% Part: incompleteness
% Chapter: incompleteness-provability
% Section: second-incompleteness-thm

\documentclass[../../../include/open-logic-section]{subfiles}

\begin{document}

\olfileid{inc}{inp}{2in}

\olsection{Teorema Ora Jangkep Kapindho}

Kepiye carane nyatakake yen $\Th{PA}$ ora mbuktekake
kekonsistenane dhewe? Nyebut $\Th{PA}$ ora konsisten padha karo
nyebut yen $\Th{PA} \Proves \eq[0][1]$. Mula kita bisa netepake
pratelan kekonsistenan $\OCon[\Th{PA}]$ minangka !!{sentence}
$\lnot \OProv[\Th{PA}](\gn{\eq[0][1]})$, lan teorema sabanjure
nyatakake asil kang dikarepake:

\begin{thm}
\ollabel{thm:second-incompleteness} 
Yen $\Th{PA}$ konsisten, mula $\Th{PA}$ ora !!{derive}
$\OCon[\Th{PA}]$.
\end{thm}

Wigati dimangerteni yen teorema iki gumantung marang representasi
tartamtu saka $\OCon[\Th{PA}]$ (yaiku, representasi tartamtu saka
$\OProv[\Th{PA}](y)$). Kang digunakake mung kasunyatan yen
representasi $\OProv[\Th{PA}](y)$ nyukupi telung sarat
!!{derivability}, mula teorema iki bisa digeneralisasi marang
saben teori kang nduweni predikat !!a{derivability} kanthi sipat-sipat
iki.

Migunani maca rancangan argumen G\"odel, amarga teorema iki banjur
dadi panutup kang ngena. Gagasane mangkene. Upamana $!G_\Th{PA}$
ukara G\"odel kang wis dibangun ing bukti
\olref[1in]{thm:first-incompleteness}. Kita wis nuduhake
``Yen $\Th{PA}$ konsisten, mula $\Th{PA}$ ora !!{derive}
$!G_\Th{PA}$.'' Yen iki diformalisasi \emph{ing} $\Th{PA}$, kita
nduweni bukti kanggo
\[
\OCon[\Th{PA}] \lif \lnot \OProv[\Th{PA}](\gn{!G_\Th{PA}}).
\]
Saiki upamana $\Th{PA}$ !!{derive}s $\OCon[\Th{PA}]$. Mula teori iki
!!{derive}s $\lnot \Prov[\Th{PA}](\gn{!G_\Th{PA}})$. Nanging amarga
$!G_\Th{PA}$ iku ukara G\"odel, pratelan iki ekuivalen karo
$!G_\Th{PA}$. Mula $\Th{PA}$ !!{derive}s $!G_\Th{PA}$. Cathetan koreksi sumber OLPL-798: ing langkah sadurunge, sumber nulis negasi Prov tanpa O, padahal argumen ing PA mbutuhake formula internal OProv kaya ing tampilan implikasi lan definisi ukara Gödel.

Nanging kita ngerti: yen $\Th{PA}$ konsisten, teori iki ora
!!{derive} $!G_\Th{PA}$!{} Dadi yen $\Th{PA}$ konsisten, teori iki
ora bisa !!{derive} $\OCon[\Th{PA}]$.

Supaya argumen iki luwih presisi, upamana $!G_\Th{PA}$ iku ukara
G\"odel kanggo~$\Th{PA}$, banjur gunakna sarat !!{derivability}
(P1)--(P3) kanggo nuduhake yen $\Th{PA}$ !!{derive}s
$\OCon[\Th{PA}] \lif !G_\Th{PA}$. Iki bakal nuduhake yen
$\Th{PA}$ ora !!{derive} $\OCon[\Th{PA}]$. Ing ngisor iki rancangan
buktine ing~$\Th{PA}$. (Supaya prasaja, subskrip $\Th{PA}$
diilangake.)
\begin{align}
& !G \liff \lnot \OProv(\gn{!G}) \ollabel{G2-1}\\
& \qquad\text{$!G$ iku ukara G\"odel}\notag \\
& !G \lif \lnot \OProv(\gn{!G}) \ollabel{G2-2}\\
  & \qquad\text{saka \olref{G2-1}} \notag\\
& !G \lif
  (\OProv(\gn{!G}) \lif \lfalse) \ollabel{G2-3}\\
  & \qquad\text{saka \olref{G2-2} miturut logika}\notag\\
& \OProv(\gn{
    !G \lif
    (\OProv(\gn{!G}) \lif \lfalse)
  }) \ollabel{G2-4}\\
  & \qquad\text{saka \olref{G2-3} miturut sarat P1} \notag\\
& \OProv(\gn{!G}) \lif
  \OProv(\gn{
    (\OProv(\gn{!G}) \lif \lfalse)
    }) \ollabel{G2-5}\\
  & \qquad\text{saka \olref{G2-4} miturut sarat P2} \notag\\
& \OProv(\gn{!G}) \lif (\OProv(\gn{\OProv(\gn{!G})}) \lif \OProv(\gn{\lfalse})) \ollabel{G2-6}\\
  & \qquad\text{saka \olref{G2-5} miturut sarat P2 lan logika} \notag\\
& \OProv(\gn{!G}) \lif 
  \OProv(\gn{\OProv(\gn{!G})}) \ollabel{G2-7}\\
   & \qquad\text{miturut P3} \notag\\
& \OProv(\gn{!G}) \lif \OProv(\gn{\lfalse}) \ollabel{G2-8}\\
  & \qquad \text{saka \olref{G2-6} lan \olref{G2-7} miturut logika}\notag\\
& \OCon \lif \lnot \OProv(\gn{!G}) \ollabel{G2-9}\\
  & \qquad\text{kontraposisi saka \olref{G2-8} lan $\OCon \ident \lnot \OProv(\gn{\lfalse})$}\notag \\
& \OCon \lif !G \notag\\
  & \qquad\text{saka \olref{G2-1} lan \olref{G2-9} miturut logika}\notag
\end{align}
Logika kang digunakake ing ndhuwur mung fakta dhasar logika
proposisional; upamane, \olref{G2-3} nggunakake
$\Proves \lnot!A \liff (!A\lif \lfalse)$ lan \olref{G2-8}
nggunakake $!A \lif (!B \lif !C), !A \lif !B \Proves !A \lif !C$.
Panganggone sarat~P2 ing \olref{G2-5} lan \olref{G2-6}
gumantung marang conto saka~P2,
$\OProv(\gn{!A \lif !B}) \lif
(\OProv(\gn{!A}) \lif \OProv(\gn{!B}))$. Ing kang kapisan,
$!A \ident !G$ lan $!B \ident \OProv(\gn{!G}) \lif \lfalse$;
ing kang kapindho, $!A \ident \OProv(\gn{G})$ lan
$!B \ident \lfalse$. Cathetan koreksi sumber OLPL-799: conto kapindho P2 kudu nggunakake numeral kode formula G kang padha karo kabeh baris G2; tandha formula ing njero numeral iku ilang ing teks sumber.

Wujud luwih abstrak saka teorema ora jangkep kapindho yaiku:

\begin{thm}
\ollabel{thm:second-incompleteness-gen} Upamana $\Th{T}$ teori
konsisten sembarang kang !!{axiomatized} lan ngembangake $\Th{Q}$,
lan upamana $\OProv[\Th{T}](y)$ formula sembarang kang nyukupi sarat
!!{derivability} P1--P3 kanggo~$\Th{T}$. Mula $\Th{T}$ ora
!!{derive}~$\OCon[T]$.
\end{thm}

\begin{prob}
Tuduhna yen $\Th{PA}$ !!{derive}s
$!G_{\Th{PA}} \lif \OCon[\Th{PA}]$.
\end{prob}

\begin{digress}
Piwulang saka crita iki yaiku ora ana teori matematika konsisten kang
``lumrah'' kang bisa !!{derive} pratelan kekonsistenane dhewe.
Upamana $\Th{T}$ teori matematika kang ngemot $\Th{Q}$ lan
panalaran ``finitari'' saka Hilbert (apa wae teges sing dimaksud).
Mula sakabehe $\Th{T}$ ora bisa !!{derive} pratelan
kekonsistenan saka $\Th{T}$, luwih-luwih bagean finitarine uga ora
bisa !!{derive} pratelan kekonsistenan saka~$\Th{T}$. Ing teges iki,
ora mungkin ana bukti kekonsistenan finitari kanggo
``sakabehe matematika.''

Ana keluwesan nalika nafsirake istilah ``finitari,'' lan ing makalah
taun 1931 G\"odel ngakoni kemungkinan yen sawijining perkara kang
kita anggep ``finitari'' bisa ana ing njaba jinis matematika
kang arep diformalisasi Hilbert. Nanging G\"odel wis menehi
kelonggaran; saiki angel mbayangake perkara kang lumrah diarani
finitari nanging ora bisa diformalisasi ing, upamane, $\Th{ZFC}$,
teori himpunan Zermelo--Fraenkel kanthi aksioma pilihan.
\end{digress}

\end{document}
